% Cible : ARIEL_LOGMINER_MEMOIRE_FINAL/frontmatter/declaration_ia.tex
% Passage à remplacer : fichier entier, de « \chapter*{Déclaration relative » à la fin
% Agent : 6 — Reproductibilité, artefacts et déclaration IA
% Sources : déclaration actuelle ; RETOURS_LECTEURS.md, lecteur 1 E.3 et lecteur 2 §2.3/point D ;
%           HANDOFF.md §3-4 ; artefacts experiments/phase_revision_20261104/ et leur manifeste.

\chapter*{Déclaration relative à l'utilisation d'outils d'intelligence artificielle}
\addcontentsline{toc}{chapter}{Déclaration relative à l'utilisation d'outils d'intelligence artificielle}

Des outils d'intelligence artificielle générative ont été employés comme assistants pour
la rédaction, la programmation, le contrôle de cohérence et l'exécution automatisée de
campagnes expérimentales. Ils ne constituent une source scientifique d'aucune affirmation
du mémoire. Les références bibliographiques renvoient aux publications, normes,
documentations et jeux de données effectivement consultés.

\section*{Outils et usages}

\begin{table}[H]
\centering
\small
\begin{tabularx}{\textwidth}{p{0.27\textwidth}X}
\toprule
\textbf{Outil} & \textbf{Usages} \\
\midrule
ChatGPT et OpenAI Codex, modèle GPT-5 (OpenAI) &
Structuration et reformulation du manuscrit ; assistance à l'écriture, à la révision et
au test de scripts ; contrôles de cohérence entre texte, code et artefacts. \\
Claude Code, modèle Claude Opus (Anthropic) &
Audit du manuscrit et du dépôt ; proposition et mise en œuvre de protocoles
complémentaires ; écriture ou adaptation de scripts ; exécution de campagnes dans
l'environnement de l'auteur ; agrégation des artefacts et préparation de corrections. \\
\bottomrule
\end{tabularx}
\end{table}

\section*{Parties concernées}

L'assistance a concerné le code du prototype et les scripts expérimentaux, la préparation
et l'exécution de campagnes complémentaires, la documentation des preuves, la génération
de figures et de tableaux à partir des artefacts, ainsi que la rédaction et
l'harmonisation du manuscrit. Les valeurs numériques retenues proviennent des fichiers
CSV ou JSON identifiés dans la matrice résultats--preuves et non de réponses textuelles
des outils.

Les campagnes complémentaires assistées par IA ne sont pas présentées comme
préenregistrées. Un protocole n'est qualifié de préenregistré que si un document
horodaté et versionné, antérieur à l'exécution, fixe les comparateurs, les métriques et
le critère de décision.

\section*{Vérification et responsabilité}

Les propositions ont été contrôlées en reliant les valeurs à un identifiant de campagne
et à un artefact, en comparant les descriptions du comportement au code et aux journaux,
et en conservant les résultats défavorables ou contradictoires. Les empreintes SHA-256
du paquet de preuves permettent de vérifier l'intégrité binaire des artefacts publiés.
Ces contrôles établissent une traçabilité ; ils ne valent pas reproduction indépendante
par un tiers.

L'auteur a arrêté les choix retenus dans le manuscrit et demeure responsable de la
problématique, des résultats présentés, de leur interprétation, du choix des références
et du contenu final. L'étendue précise de l'assistance est déclarée ci-dessus afin que
les apports automatisés ne soient pas confondus avec une contribution humaine non
assistée.
